\chapter{Аналитическая часть}
В данной лабораторной работе рассматривается 2 алгоритма поиск элемента в массиве: 
\begin{itemize}
	\item поиск полным перебором
	\item бинарный поиск
\end{itemize}

\section{Поиск полным перебором}
Пусть дан массив  из N различных элементов. Рассматривается каждый элемент массива и сравнивается с искомым. Если рассматриваемый элемент совпадает с искомым его индекс запоминается и считается результатом поиск. Общая сложность алгоритма $O(N)$.

\section{Бинарный поиск}
Пусть дан отсортированный по возрастанию массив $array$ из N различных чисел, искомый элемент $t$. $l=0, r=N-1$ текущие индексы левой и правой границы массива соответственно. В каждой итерации поиска определяется средний элемент массива с индексом 

\begin{equation}
	\label{eq:bin_s_m}
	m_i = l + \frac{r - l}{2}
\end{equation}

Для следующей итерации определяются новый границы массива
\begin{equation}
	\label{eq:new_l}
	l = 
	\begin{cases}
		m+1 & \text{array[m] < t} \\
		l & \text{иначе}
	\end{cases}
\end{equation}
\begin{equation}
	\label{eq:new_r}
	r = 
	\begin{cases}
		m-1 & \text{array[m]} \geq \text{t} \\
		r & \text{иначе}
	\end{cases}
\end{equation}
Поиск продолжается до тех пор пока $l < r$

Уменьшение размера массива в 2 раза на каждой итерации позволяет добавить временной сложности  $O(\log{N})$, но при этом требуется чтобы массив поступал и поддерживался отсортированным

\section{Вывод}
Алгоритм бинарного поиска является более более оптимизированный методом поиска элемента в массиве, но накладывает дополнительные ограничения --- поддержка массива отсортированным. При этом алгоритм полного перебора требует больше временных затрат, но при этом нет требований для входного массива
